% Original: PDF strana 26, donji slajd.
\slajdid{02-s052}
\begin{frame}[label=02-s052]{Minimizacija i prelaz između površina}
% element: 02-s052:e01
\Rezultat{\(\textcolor{DEcrvena}{F=\bar C\bar B+BA}\)}
\par\vspace{4mm}
% element: 02-s052:e02
\begin{columns}[T,onlytextwidth]
\begin{column}{.30\textwidth}\centering
\Oznaka{Član \(\bar C\bar B\)}\vspace{4mm}
% Izvorne karte su očuvane; formula je usklađena po korisnikovoj odluci.
\begingroup
\fontsize{16.364}{18.2}\selectfont
\renewcommand{\small}{\fontsize{16.364}{18.2}\selectfont}
\scalebox{.55}{%
\begin{karnaugh-map}(label=middle,implicantcolors={DEcrvena,DEcrvena,DEcrvena})[4][2][1][$A$][$B$][$C$]
\terms{0,1,2,3,4,5,6,7}{\strut}
\implicant{0}{1}
\end{karnaugh-map}}
\endgroup

\end{column}
\begin{column}{.30\textwidth}\centering
\Oznaka{Član \(BA\)}\vspace{4mm}
% Izvorne karte su očuvane; formula je usklađena po korisnikovoj odluci.
\begingroup
\fontsize{16.364}{18.2}\selectfont
\renewcommand{\small}{\fontsize{16.364}{18.2}\selectfont}
\scalebox{.55}{%
\begin{karnaugh-map}(label=middle,implicantcolors={DEcrvena,DEcrvena,DEcrvena})[4][2][1][$A$][$B$][$C$]
\terms{0,1,2,3,4,5,6,7}{\strut}
\implicant{3}{7}
\end{karnaugh-map}}
\endgroup

\end{column}
\begin{column}{.34\textwidth}\centering
\Oznaka{Karta funkcije}\vspace{4mm}
% Izvorne karte su očuvane; formula je usklađena po korisnikovoj odluci.
\begingroup
\fontsize{16.364}{18.2}\selectfont
\renewcommand{\small}{\fontsize{16.364}{18.2}\selectfont}
\scalebox{.55}{%
\begin{karnaugh-map}(label=middle,implicantcolors={DEcrvena,DEcrvena,DEcrvena})[4][2][1][$A$][$B$][$C$]
\minterms{0,1,3,7}\maxterms{2,4,5,6}
\implicant{0}{1}
\end{karnaugh-map}}
\endgroup

\end{column}
\end{columns}
\par\vspace{5mm}
Za promenljive \(C,B,A\), funkcija nastaje minimizacijom potpunog oblika prikazanog u Karnoovoj karti.
\end{frame}
